\nwfilename{nw/unification.nw}\nwbegindocs{0}% -*- mode: poly-noweb; noweb-code-mode: sml-mode; -*-% ===> this file was generated automatically by noweave --- better not edit it
%%
%% unification.nw
%% 
%% Made by Alex Nelson <pqnelson@gmail.com>
%% Login   <alex@lisp>
%% 
%% Started on  2026-04-25T15:53:23-0700
%% Last update 2026-04-25T15:53:23-0700
%% 

\section{Unification}


\begin{node}
Unification, for first-order logic, given two ``expressions'' $e_{1}$
and $e_{2}$ (a pair of terms or a pair of formulas) tries to find a
``substitution'' $\sigma$ which can be applied to obtain
\emph{syntactically} identical expressions
$\sigma(e_{1})=\sigma(e_{2})$ (where $\sigma(-)$ essentially
corresponds to a ``parallel'' version of {\Tt{}Formula.subst\nwendquote} or {\Tt{}Term.subst\nwendquote}). The {\Tt{}Formula.eq\nwendquote} and {\Tt{}Term.eq\nwendquote}
predicates test if two formulas (resp., terms) are syntactically
identical. If two expressions are unifiable, then it returns such a
substitution. 

We will need only ``first-order unification'', i.e., we work with only
``first-order variables'' and not ``higher-order
variables''. Higher-order logic allows us to quantify over predicates
and functions using higher-order variables, but unification for
higher-order variables is rather complicated. First-order unification
merely allows replacing variables with terms, which is far more tame
by comparison.

We will only need to unify terms and literals (predicates or their
negation). 
\end{node}

\nwenddocs{}\nwbegincode{1}\sublabel{NW1BOcHB-3z172B-1}\nwmargintag{{\nwtagstyle{}\subpageref{NW1BOcHB-3z172B-1}}}\moddef{Unif.sig~{\nwtagstyle{}\subpageref{NW1BOcHB-3z172B-1}}}\endmoddef\nwstartdeflinemarkup\nwenddeflinemarkup
signature Unif = sig
  structure Env :
            sig
              type t;
              val undefined : unit -> t;
              val apply : t -> Term.t -> Term.t;
              val defined : t -> Term.t -> bool;
              val insert : t -> Term.t -> Term.t -> t;
            end;
  val y : \nwlinkedidentc{Env.t}{NW1BOcHB-25AKca-1} -> (Term.t * Term.t) list -> \nwlinkedidentc{Env.t}{NW1BOcHB-25AKca-1};
  val y_literals : \nwlinkedidentc{Env.t}{NW1BOcHB-25AKca-1} -> Formula.t * Formula.t -> \nwlinkedidentc{Env.t}{NW1BOcHB-25AKca-1};
  val y_complements : \nwlinkedidentc{Env.t}{NW1BOcHB-25AKca-1} -> Formula.t * Formula.t -> \nwlinkedidentc{Env.t}{NW1BOcHB-25AKca-1};
end;
\nwnotused{Unif.sig}\nwidentuses{\\{{\nwixident{Env.t}}{Env.t}}\\{{\nwixident{Formula.t}}{Formula.t}}\\{{\nwixident{Term.t}}{Term.t}}}\nwindexuse{\nwixident{Env.t}}{Env.t}{NW1BOcHB-3z172B-1}\nwindexuse{\nwixident{Formula.t}}{Formula.t}{NW1BOcHB-3z172B-1}\nwindexuse{\nwixident{Term.t}}{Term.t}{NW1BOcHB-3z172B-1}\nwendcode{}\nwbegindocs{2}\nwdocspar

\nwenddocs{}\nwbegincode{3}\sublabel{NW1BOcHB-1KP4Zf-1}\nwmargintag{{\nwtagstyle{}\subpageref{NW1BOcHB-1KP4Zf-1}}}\moddef{Unif.sml~{\nwtagstyle{}\subpageref{NW1BOcHB-1KP4Zf-1}}}\endmoddef\nwstartdeflinemarkup\nwenddeflinemarkup
structure Unif :> Unif = struct
\LA{}Finite map for unification~{\nwtagstyle{}\subpageref{NW1BOcHB-25AKca-1}}\RA{}

\LA{}Poor man's hashmap~{\nwtagstyle{}\subpageref{nw@notdef}}\RA{}

\LA{}Unify two terms~{\nwtagstyle{}\subpageref{NW1BOcHB-1kcz6k-1}}\RA{}

\LA{}Unify two literals~{\nwtagstyle{}\subpageref{NW1BOcHB-3w0JS3-1}}\RA{}
end;
\nwnotused{Unif.sml}\nwendcode{}\nwbegindocs{4}\nwdocspar

\begin{node}
Unification tracks the substitution constructed using a finite partial
map. We recursively build the map from equations given for unification.
We will write $t\sim t'$ for the ``equation'' submitted for
unification involving the two terms $t$ and $t'$. This will produce a
finite partial map, and we will denote a specific entry in it as
$x\mapsto t$ to indicate $x$ would be substituted for the term $t$.
\end{node}

\begin{definition}
Suppose $M$ is a finite partial map from terms to terms.
\begin{enumerate}
\item Let $x\mapsto t$ be an entry in the map $M$. We write $x\to y$ if there exists a $y\in FV(t)$.
\item A \define{Cycle} in $M$ is a nonempty finite sequence leading
  back to the starting point $x_{0}\to x_{1}\to\dots\to x_{n}\to x_{0}$.
\item If there are no cycles in $M$, we say it is \define{Cycle-free}. Otherwise,
  we call $M$ \define{Cyclic}.
\end{enumerate}
\end{definition}

\begin{example}[{Duffy~\cite[p.143]{duffy1991principles}}]
Consider unifying $P[x,f(x)]$ with $P[g(y),y]$ where $P[-,-]$ is a
predicate of two terms, $f(-)$ and $g(-)$ are functions.

The first step will be to break this up into two equations
\begin{equation}
x\sim g(y)
\end{equation}
and
\begin{equation}
f(x)\sim y.
\end{equation}
We can solve the first equation with the replacement $x\mapsto g(y)$.

The second equation can be ``rewritten'' as
\begin{equation}
y\sim f(x),
\end{equation}
by symmetry. However, if we try adjoining the replacement $y\mapsto f(x)$
to our current partial map, then we end up with a cycle $x\to y\to x$.

We want to avoid this. Why? Because if we consider the formula,
\begin{equation}\label{eq:unification:example-of-why-we-want-to-avoid-cycles}
(\forall x\exists y\ldotp P[x,y])\implies(\exists y\forall x\ldotp P[x,y]),
\end{equation}
the negation of this formula has conjunctive normal form
\begin{equation}
\forall x,y\ldotp P[x,f(x)]\land\neg P[g(y),y].
\end{equation}
A resolution proof can be found for this, if we allowed cycles during
unification. But Equation~\eqref{eq:unification:example-of-why-we-want-to-avoid-cycles}
is not a theorem in classical logic. Had we allowed cycles, soundness
would be compromised. 
\end{example}

\begin{node}
The finite map used for tracking variable assignment could be
anything. For simplicity, at the moment, we use association lists.

A ``Poor Man's Table'' is just an association list.
If we implement a more optimized hash table, like a Red-Black tree or
a 2-3 tree, then we can jettison this implementation and use something
superior. The contract:
\begin{enumerate}
\item {\Tt{}undefined\nwendquote} is an empty table
\item {\Tt{}apply\ table\ key\nwendquote} gets the value associated with the {\Tt{}key\nwendquote}
  in the {\Tt{}table\nwendquote}
\item {\Tt{}insert\ table\ key\ value\nwendquote} inserts the association of {\Tt{}value\nwendquote}
  with {\Tt{}key\nwendquote} into the {\Tt{}table\nwendquote}, raising a {\Tt{}Fail\nwendquote} exception if
  this would produce a cycle in the substitution;
\item {\Tt{}defined\ table\ key\nwendquote} tests if the {\Tt{}table\nwendquote} contains the
  {\Tt{}key\nwendquote} at all
\end{enumerate}
These are the table operations.

Right now, this is a huge bottleneck for performance. The low-hanging
fruit would be to make this replacement.
\end{node}

\begin{invariant}[Partial maps remain cycle free]\label{invariant:unification:cycle-free}
Every [public] function in the {\Tt{}Unif\nwendquote} structure requires any
{\Tt{}\nwlinkedidentq{Env.t}{NW1BOcHB-25AKca-1}\nwendquote} partial map given as an argument to be cycle free, and
ensures the {\Tt{}\nwlinkedidentq{Env.t}{NW1BOcHB-25AKca-1}\nwendquote} partial map produced (if any) to be cycle-free.
\end{invariant}

\nwenddocs{}\nwbegincode{5}\sublabel{NW1BOcHB-25AKca-1}\nwmargintag{{\nwtagstyle{}\subpageref{NW1BOcHB-25AKca-1}}}\moddef{Finite map for unification~{\nwtagstyle{}\subpageref{NW1BOcHB-25AKca-1}}}\endmoddef\nwstartdeflinemarkup\nwusesondefline{\\{NW1BOcHB-1KP4Zf-1}}\nwenddeflinemarkup
structure Env :>
            sig
              type t;
              val undefined : unit -> t;
              val apply : t -> Term.t -> Term.t;
              val defined : t -> Term.t -> bool;
              val insert : t -> Term.t -> Term.t -> t;
            end =
struct
type t = (Term.t * Term.t) list;

fun undefined () = [];

fun apply [] _ = raise Fail "apply"
  | apply ((y,tm)::ys) x = if Term.eq x y
                           then tm
                           else apply ys x;

fun defined [] _ = false
  | defined ((y,_)::env) x = (Term.eq x y) orelse
                             defined env x;
local
\LA{}Test if a new equations is trivial for unification~{\nwtagstyle{}\subpageref{NW1BOcHB-4edp1Q-1}}\RA{}
in
fun insert env (x as (Term.Var _)) y = 
      if is_cyclic env x y then env 
      else (x,y)::env
  | insert _ _ _ = raise Fail "\nwlinkedidentc{Env.insert}{NW1BOcHB-25AKca-1}: trying to insert function as key"
end;
end;
\nwindexdefn{\nwixident{Env.t}}{Env.t}{NW1BOcHB-25AKca-1}\nwindexdefn{\nwixident{Env.undefined}}{Env.undefined}{NW1BOcHB-25AKca-1}\nwindexdefn{\nwixident{Env.apply}}{Env.apply}{NW1BOcHB-25AKca-1}\nwindexdefn{\nwixident{Env.insert}}{Env.insert}{NW1BOcHB-25AKca-1}\nwindexdefn{\nwixident{Env.defined}}{Env.defined}{NW1BOcHB-25AKca-1}\eatline
\nwused{\\{NW1BOcHB-1KP4Zf-1}}\nwidentdefs{\\{{\nwixident{Env.apply}}{Env.apply}}\\{{\nwixident{Env.defined}}{Env.defined}}\\{{\nwixident{Env.insert}}{Env.insert}}\\{{\nwixident{Env.t}}{Env.t}}\\{{\nwixident{Env.undefined}}{Env.undefined}}}\nwidentuses{\\{{\nwixident{Term.eq}}{Term.eq}}\\{{\nwixident{Term.t}}{Term.t}}}\nwindexuse{\nwixident{Term.eq}}{Term.eq}{NW1BOcHB-25AKca-1}\nwindexuse{\nwixident{Term.t}}{Term.t}{NW1BOcHB-25AKca-1}\nwendcode{}\nwbegindocs{6}\nwdocspar
\begin{node}
We also have the really interesting part of solving the unification
problem: testing if $x=t$ is a trivial equation or not. This would
happen for literally trivial equations $x=x$ which shed no light on anything,
but also for $x=y$ when we already have $y=t'$.

However, there is the possibility that we might produce a cycle. This
is handled by raising an error.

Specifically: checking there are no cycles when we try to associate
variable $x$ with term $t$ depends on what $t$ is. When
\begin{enumerate}
\item $t$ is a variable $t=y$: This boils down to two subcases.
  \begin{enumerate}
  \item $x=y$ 
  \item $x\neq y$, but $y$ is bound to something in the
    environment $y=t_{y}$. Then we just check recursively that associating $x$ to
    $t_{y}$ will be cycle-free.
  \end{enumerate}
\item $t$ is a function $f(\vec{t})$: There is no variable $y\in
  FV(t)$ such that $y\to^{*}x$ (i.e., there is no sequence of zero or more
  steps leading from $y$ to $x$)
\end{enumerate}
We can check the first condition elsewhere (namely, in the {\Tt{}unify\nwendquote}
function). The second condition will be tested in {\Tt{}is{\_}cyclic\nwendquote},
which will return \textbf{false} if the condition is satisfied.
\end{node}

\begin{node}
If we extend {\Tt{}env\nwendquote} with a new entry $x\sim t$, then we should
consider the following situations:
\begin{enumerate}
\item $t$ is a variable $t=y$.
  \begin{enumerate}
  \item If $y=x$ are the same variable, then this is a trivial
    extension, and we should return \textbf{true} to indicate \emph{do
    no extend the environment}.
  \item If $y\neq x$ is a different variable, then we should check if
    $y\sim t'$ is in the environment and if so whether we can add
    $x\sim t'$ to the environment.

    If there is no entry $y\sim t'$ in the environment, then return
    \textbf{false} to indicate \emph{we can extend the environment
    while maintaining the cycle-free invariant}.

    If there is an entry $y\sim t'$ in the environment, then we just
    recursively test if $x\sim t'$ violates the cycle-free invariant
    of the environment.
  \end{enumerate}
\item If $t$ is a function $f(t_{1},\dots,t_{n})$, then we want to test 
  that none of the possible extensions $x\sim t_{1}$, \dots, $x\sim t_{n}$
  would lead to a cyclic environment. If there is a possible extension
  $x\sim t_{k}$, then return \textbf{true}. Otherwise,
  return \textbf{false} to indicate \emph{we can extend the
  environment while maintaining the cycle-free invariant}.
\end{enumerate}
\end{node}

\nwenddocs{}\nwbegincode{7}\sublabel{NW1BOcHB-4edp1Q-1}\nwmargintag{{\nwtagstyle{}\subpageref{NW1BOcHB-4edp1Q-1}}}\moddef{Test if a new equations is trivial for unification~{\nwtagstyle{}\subpageref{NW1BOcHB-4edp1Q-1}}}\endmoddef\nwstartdeflinemarkup\nwusesondefline{\\{NW1BOcHB-25AKca-1}}\nwenddeflinemarkup
(* REQUIRES: `x` is not defined in the `env` *)
fun is_cyclic env x tm =
 (case tm of
   (y as (Term.Var _)) => Term.eq x y orelse (* trivial case *)
                          Term.have_same_sort x y andalso
                          defined env y andalso
                          is_cyclic env x (apply env y)
 | (Term.Fun(f,s,args)) => List.exists (is_cyclic env x) args andalso
                           raise Fail "cyclic");
\nwused{\\{NW1BOcHB-25AKca-1}}\nwidentuses{\\{{\nwixident{Term.eq}}{Term.eq}}\\{{\nwixident{Term.have{\_}same{\_}sort}}{Term.have:unsame:unsort}}}\nwindexuse{\nwixident{Term.eq}}{Term.eq}{NW1BOcHB-4edp1Q-1}\nwindexuse{\nwixident{Term.have{\_}same{\_}sort}}{Term.have:unsame:unsort}{NW1BOcHB-4edp1Q-1}\nwendcode{}\nwbegindocs{8}\nwdocspar

\begin{node}
Unification tries to ``solve'' a bunch of equations. We have the
following situations:
\begin{enumerate}
\item $f(t_{1},\dots,t_{m})=g(t'_{1},\dots,t'_{n})$ requires $f$ and
  $g$ to be the same function symbol with the same sort, and $m=n$,
  then we just need to solve the equations $t_{1}=t'_{1}$, \dots, $t_{n}=t'_{n}$.
\item $x=t$ first requires checking if we have already ``solved'' for
  $x$, and if so try to match it against $t$.

  Otherwise, if $x=t$ is a trivial equation, then we just move along
  to try solving the remaining equations.

  If $x=t$ is not a trivial equation, then we extend the current
  environment (``solutions'') with this entry, and then try solving
  the remaining equations.
\item $t=x$ is just the same situation.
\end{enumerate}
There cannot be any other possible situation.
\end{node}

\nwenddocs{}\nwbegincode{9}\sublabel{NW1BOcHB-1kcz6k-1}\nwmargintag{{\nwtagstyle{}\subpageref{NW1BOcHB-1kcz6k-1}}}\moddef{Unify two terms~{\nwtagstyle{}\subpageref{NW1BOcHB-1kcz6k-1}}}\endmoddef\nwstartdeflinemarkup\nwusesondefline{\\{NW1BOcHB-1KP4Zf-1}}\nwenddeflinemarkup
(* \nwlinkedidentc{Unif.y}{NW1BOcHB-1kcz6k-1} : \nwlinkedidentc{Env.t}{NW1BOcHB-25AKca-1} -> (Term.t * Term.t) list -> \nwlinkedidentc{Env.t}{NW1BOcHB-25AKca-1}

REQUIRES: `env` is cycle-free
REQUIRES: the set `env @ eqns` has exactly the same set of
          unifiers as the original problem. *)
fun y env ([] : (Term.t * Term.t) list) = env
  | y env ((Term.Fun(f,s1,args1),Term.Fun(g,s2,args2))::eqns) = 
      if f = g andalso s1 = s2 andalso length args1 = length args2
      then y env ((ListPair.zip(args1, args2)) @ eqns)
      else raise Fail "impossible unification"
  | y env (((x as (Term.Var _)), t)::eqns) =
      if not(Term.have_same_sort x t)
      then raise Fail "impossible unification: variable has different sort than term"
      else if \nwlinkedidentc{Env.defined}{NW1BOcHB-25AKca-1} env x (* x already in env *)
      then y env ((\nwlinkedidentc{Env.apply}{NW1BOcHB-25AKca-1} env x, t)::eqns)
      else y (\nwlinkedidentc{Env.insert}{NW1BOcHB-25AKca-1} env x t) eqns
  | y env ((t, x as (Term.Var _))::eqns) = y env ((x,t)::eqns);
\nwindexdefn{\nwixident{Unif.y}}{Unif.y}{NW1BOcHB-1kcz6k-1}\eatline
\nwused{\\{NW1BOcHB-1KP4Zf-1}}\nwidentdefs{\\{{\nwixident{Unif.y}}{Unif.y}}}\nwidentuses{\\{{\nwixident{Env.apply}}{Env.apply}}\\{{\nwixident{Env.defined}}{Env.defined}}\\{{\nwixident{Env.insert}}{Env.insert}}\\{{\nwixident{Env.t}}{Env.t}}\\{{\nwixident{Term.have{\_}same{\_}sort}}{Term.have:unsame:unsort}}\\{{\nwixident{Term.t}}{Term.t}}}\nwindexuse{\nwixident{Env.apply}}{Env.apply}{NW1BOcHB-1kcz6k-1}\nwindexuse{\nwixident{Env.defined}}{Env.defined}{NW1BOcHB-1kcz6k-1}\nwindexuse{\nwixident{Env.insert}}{Env.insert}{NW1BOcHB-1kcz6k-1}\nwindexuse{\nwixident{Env.t}}{Env.t}{NW1BOcHB-1kcz6k-1}\nwindexuse{\nwixident{Term.have{\_}same{\_}sort}}{Term.have:unsame:unsort}{NW1BOcHB-1kcz6k-1}\nwindexuse{\nwixident{Term.t}}{Term.t}{NW1BOcHB-1kcz6k-1}\nwendcode{}\nwbegindocs{10}\nwdocspar
\begin{node}
Unifying two literals $P[t_{1},\dots,t_{m}]$ with $Q[u_{1},\dots,u_{n}]$
requires checking $P=Q$ (they have the same identifier/name) and $m=n$ (they have the same number of arguments), then invoking
{\Tt{}unify\nwendquote} to solve the system $t_{1}=u_{1}$, \dots, $t_{n}=u_{n}$.
If $P\neq Q$ are distinct identifiers, or $m\neq n$, then a failure is raised.
If $Q$ is the negation of $P$, then a failure is raised (there is no
way to unify a predicate with its negation).
\end{node}

\nwenddocs{}\nwbegincode{11}\sublabel{NW1BOcHB-3w0JS3-1}\nwmargintag{{\nwtagstyle{}\subpageref{NW1BOcHB-3w0JS3-1}}}\moddef{Unify two literals~{\nwtagstyle{}\subpageref{NW1BOcHB-3w0JS3-1}}}\endmoddef\nwstartdeflinemarkup\nwusesondefline{\\{NW1BOcHB-1KP4Zf-1}}\nwenddeflinemarkup
(* \nwlinkedidentc{Unif.y_literals}{NW1BOcHB-3w0JS3-1} : \nwlinkedidentc{Env.t}{NW1BOcHB-25AKca-1} -> Term.t * Term.t -> \nwlinkedidentc{Env.t}{NW1BOcHB-25AKca-1} *)
fun y_literals env (lhs,rhs) =
  if Formula.is_not lhs andalso Formula.is_not rhs
  then y_literals env (Formula.dest_not lhs, Formula.dest_not rhs)
  else if Formula.is_falsum lhs andalso Formula.is_falsum rhs
  then env
  else if Formula.is_pred lhs andalso Formula.is_pred rhs
  then let val (p1,args1) = Formula.dest_pred lhs;
           val (p2,args2) = Formula.dest_pred rhs;
       in if p1 = p2 andalso length args1 = length args2
          then y env (ListPair.zip(args1, args2))
          else raise Fail "\nwlinkedidentc{Unif.y_literals}{NW1BOcHB-3w0JS3-1}: impossible unification"
       end
  else raise Fail "\nwlinkedidentc{Unif.y_literals}{NW1BOcHB-3w0JS3-1}: Cannot unify literals";

fun negate A = if Formula.is_not A
               then Formula.dest_not A
               else Formula.mk_not A;

fun y_complements env (A,B) =
  y_literals env (A, negate B);
\nwindexdefn{\nwixident{Unif.y{\_}complements}}{Unif.y:uncomplements}{NW1BOcHB-3w0JS3-1}\nwindexdefn{\nwixident{Unif.y{\_}literals}}{Unif.y:unliterals}{NW1BOcHB-3w0JS3-1}\eatline
\nwused{\\{NW1BOcHB-1KP4Zf-1}}\nwidentdefs{\\{{\nwixident{Unif.y{\_}complements}}{Unif.y:uncomplements}}\\{{\nwixident{Unif.y{\_}literals}}{Unif.y:unliterals}}}\nwidentuses{\\{{\nwixident{Env.t}}{Env.t}}\\{{\nwixident{Formula.dest{\_}not}}{Formula.dest:unnot}}\\{{\nwixident{Formula.dest{\_}pred}}{Formula.dest:unpred}}\\{{\nwixident{Formula.is{\_}falsum}}{Formula.is:unfalsum}}\\{{\nwixident{Formula.is{\_}not}}{Formula.is:unnot}}\\{{\nwixident{Formula.is{\_}pred}}{Formula.is:unpred}}\\{{\nwixident{Formula.mk{\_}not}}{Formula.mk:unnot}}\\{{\nwixident{Term.t}}{Term.t}}}\nwindexuse{\nwixident{Env.t}}{Env.t}{NW1BOcHB-3w0JS3-1}\nwindexuse{\nwixident{Formula.dest{\_}not}}{Formula.dest:unnot}{NW1BOcHB-3w0JS3-1}\nwindexuse{\nwixident{Formula.dest{\_}pred}}{Formula.dest:unpred}{NW1BOcHB-3w0JS3-1}\nwindexuse{\nwixident{Formula.is{\_}falsum}}{Formula.is:unfalsum}{NW1BOcHB-3w0JS3-1}\nwindexuse{\nwixident{Formula.is{\_}not}}{Formula.is:unnot}{NW1BOcHB-3w0JS3-1}\nwindexuse{\nwixident{Formula.is{\_}pred}}{Formula.is:unpred}{NW1BOcHB-3w0JS3-1}\nwindexuse{\nwixident{Formula.mk{\_}not}}{Formula.mk:unnot}{NW1BOcHB-3w0JS3-1}\nwindexuse{\nwixident{Term.t}}{Term.t}{NW1BOcHB-3w0JS3-1}\nwendcode{}

\nwixlogsorted{c}{{Finite map for unification}{NW1BOcHB-25AKca-1}{\nwixu{NW1BOcHB-1KP4Zf-1}\nwixd{NW1BOcHB-25AKca-1}}}%
\nwixlogsorted{c}{{Poor man's hashmap}{nw@notdef}{\nwixu{NW1BOcHB-1KP4Zf-1}}}%
\nwixlogsorted{c}{{Test if a new equations is trivial for unification}{NW1BOcHB-4edp1Q-1}{\nwixu{NW1BOcHB-25AKca-1}\nwixd{NW1BOcHB-4edp1Q-1}}}%
\nwixlogsorted{c}{{Unif.sig}{NW1BOcHB-3z172B-1}{\nwixd{NW1BOcHB-3z172B-1}}}%
\nwixlogsorted{c}{{Unif.sml}{NW1BOcHB-1KP4Zf-1}{\nwixd{NW1BOcHB-1KP4Zf-1}}}%
\nwixlogsorted{c}{{Unify two literals}{NW1BOcHB-3w0JS3-1}{\nwixu{NW1BOcHB-1KP4Zf-1}\nwixd{NW1BOcHB-3w0JS3-1}}}%
\nwixlogsorted{c}{{Unify two terms}{NW1BOcHB-1kcz6k-1}{\nwixu{NW1BOcHB-1KP4Zf-1}\nwixd{NW1BOcHB-1kcz6k-1}}}%
\nwixlogsorted{i}{{\nwixident{Env.apply}}{Env.apply}}%
\nwixlogsorted{i}{{\nwixident{Env.defined}}{Env.defined}}%
\nwixlogsorted{i}{{\nwixident{Env.insert}}{Env.insert}}%
\nwixlogsorted{i}{{\nwixident{Env.t}}{Env.t}}%
\nwixlogsorted{i}{{\nwixident{Env.undefined}}{Env.undefined}}%
\nwixlogsorted{i}{{\nwixident{Formula.dest{\_}not}}{Formula.dest:unnot}}%
\nwixlogsorted{i}{{\nwixident{Formula.dest{\_}pred}}{Formula.dest:unpred}}%
\nwixlogsorted{i}{{\nwixident{Formula.eq}}{Formula.eq}}%
\nwixlogsorted{i}{{\nwixident{Formula.is{\_}falsum}}{Formula.is:unfalsum}}%
\nwixlogsorted{i}{{\nwixident{Formula.is{\_}not}}{Formula.is:unnot}}%
\nwixlogsorted{i}{{\nwixident{Formula.is{\_}pred}}{Formula.is:unpred}}%
\nwixlogsorted{i}{{\nwixident{Formula.mk{\_}not}}{Formula.mk:unnot}}%
\nwixlogsorted{i}{{\nwixident{Formula.subst}}{Formula.subst}}%
\nwixlogsorted{i}{{\nwixident{Formula.t}}{Formula.t}}%
\nwixlogsorted{i}{{\nwixident{Term.eq}}{Term.eq}}%
\nwixlogsorted{i}{{\nwixident{Term.have{\_}same{\_}sort}}{Term.have:unsame:unsort}}%
\nwixlogsorted{i}{{\nwixident{Term.subst}}{Term.subst}}%
\nwixlogsorted{i}{{\nwixident{Term.t}}{Term.t}}%
\nwixlogsorted{i}{{\nwixident{Unif.y}}{Unif.y}}%
\nwixlogsorted{i}{{\nwixident{Unif.y{\_}complements}}{Unif.y:uncomplements}}%
\nwixlogsorted{i}{{\nwixident{Unif.y{\_}literals}}{Unif.y:unliterals}}%

